% ===================================================================
% SKILL: SUNUM BAŞLANGIÇ VE KAPAK ŞABLONU
% ===================================================================
% TANIM:
% Bu dosya, tüm hafta sunumlarında standart olarak kullanılması gereken:
% 1. Meta Veri Tanımlarını (Ders, Hoca, Tarih vb.)
% 2. Özel Görsel Kapak Tasarımını (Logo ve Arka Plan)
% 3. Otomatik İçindekiler Yapısını
% içerir. Yeni bir sunuma başlarken bu yapıyı kopyalayarak başlayın.
% ===================================================================

% -------------------------------------------------------------------
% 1. PREAMBLE: BELGE VE PAKET TANIMLARI
% -------------------------------------------------------------------
% Not: Ana .tex dosyanızın en başında bu sınıf ve paketler olmalıdır.
% \documentclass[light]{lutbeamer}
\usepackage{pgfpages}
\usepackage{ragged2e}
\usepackage{booktabs}
\usepackage{amssymb}
\usepackage{fontawesome5}
\usepackage{tikz}
\usepackage{pgfplots}
\pgfplotsset{compat=1.18}
\addbibresource{refs.bib} % BibLaTeX Kaynakça Dosyası

% -------------------------------------------------------------------
% 2. METADATA (DERS KİMLİK BİLGİLERİ)
% -------------------------------------------------------------------
% Bu kısım \begin{document} öncesinde (preamble) yer almalıdır.

\setdepartment{OSTİM TEKNİK ÜNİVERSİTESİ}                   % 1. Kurum Adı
\institute[Yazılım Mühendisliği]{YAZILIM MÜHENDİSLİĞİ BÖLÜMÜ} % 2. Bölüm Adı
\author[Mühendislik Fakültesi]{\\Prof.~Dr.~Yalçın ATA,\\ 
Prof.~Dr.~Arif DEMİR,\\ Dr.~Öğr.~Üyesi~Muhammed ELMNEFİ,\\ Dr.~Öğr.~Üyesi~Haydar KILIÇ,\\ Dr.~Öğr.~Üyesi~Emel GÜVEN,\\ Dr.~Öğr.~Üyesi~Ufuk ASIL,\\ Öğr.~Gör.~Sema ÇİFTÇİ} % 3. Eğitmen Adı

\title{MATH 204 -- Olasılık ve İstatistik}                     % 4. Ders Kodu ve Adı
% \subtitle değişkeni her hafta değişecektir:
\subtitle{Hafta X\@: [Konu Başlığı Buraya]}                   % 5. Hafta Konusu

\date{2025--2026 Bahar Dönemi}                                           % 6. Dönem

% -------------------------------------------------------------------
% 3. OTOMATİK İÇİNDEKİLER (SUNUM AKIŞI)
% -------------------------------------------------------------------
% Her \section{} komutundan önce otomatik bir "Sunum Akışı" slaytı ekler.
% Sol ve sağ sütun olarak ikiye böler (Uzun içerikler için idealdir).

\AtBeginSection[]{
  \begin{frame}{Sunum Akışı}
  \footnotesize
  \addtobeamertemplate{section in toc}{}{\vspace{0.15cm}}
  \begin{columns}[t]
    \column{.5\textwidth}
    \tableofcontents[sections={1--3},currentsection] % İlk 3 bölüm solda
    \column{.5\textwidth}
    \tableofcontents[sections={4--10},currentsection] % Kalan bölümler sağda
  \end{columns}
  \end{frame}
}

% -------------------------------------------------------------------
% 4. KAPAK SLAYTLARI (DOCUMENT BODY İÇİ)
% -------------------------------------------------------------------
% Aşağıdaki kod bloğu \begin{document} komutundan HEMEN SONRA gelmelidir.

% Adım A: Navigasyonu Gizle (Kapakta çirkin durmaması için)
\setbeamertemplate{navigation symbols}{}

% Adım B: Görsel Tam Kapak (Logolu)
\begin{frame}<article:0>[plain,noframenumbering]
    \begin{tikzpicture}[remember picture,overlay]
        % Arka Plan Görseli
        \node[anchor=center, at=(current page.center), inner sep=0pt] {
            \includegraphics[width=\paperwidth, height=\paperheight]{logos/background_figure.jpg}
        };
        % Üniversite Logosu (Ortalanmış)
        \node[anchor=center, at=(current page.center), inner sep=0pt] {
            \includegraphics[keepaspectratio, width=0.65\paperwidth, height=\paperheight]{logos/LUT-LOGO-WHITE-PNG.png}
        };
    \end{tikzpicture}
\end{frame}

% Adım C: Resmi Başlık Sayfası (Ders Adı ve Hoca Bilgileri)
\maketitle{}
